\chapter{Anesthesia}

\section*{Definition and Overview}
Anesthesia is the use of medicines to prevent pain during surgery and other procedures. It works by blocking or interrupting nerve signals so that a person feels no pain, and it may also include keeping the person asleep, relaxing muscles, and controlling breathing.

There are three main types: general anesthesia, which puts the whole body asleep; regional anesthesia, which numbs a large area such as the lower body; and local anesthesia, which numbs a small area. Anesthesia is delivered and monitored by a specialist doctor called an anesthesiologist or anaesthetist, who tailors the technique to the person and the procedure.

\section*{Epidemiology}
Anesthesia is one of the most commonly performed medical procedures: millions of operations and procedures requiring anesthesia are done every year, and it is considered very safe. Serious complications are rare, but the risk varies with the person's age, general health, and the urgency and type of surgery.

Very young and very old people, and those with heart, lung, or kidney disease, obesity, or sleep apnea, have higher risks. Most healthy people having planned surgery will recover from anesthesia with no significant problems.

\section*{Aetiology and Pathophysiology}
\begin{itemize}
    \item \textbf{General anesthesia:} inhaled gases or intravenous medicines depress the central nervous system, producing unconsciousness, pain relief, and muscle relaxation
    \item \textbf{Regional anesthesia:} an anesthetic is injected near a group of nerves, such as in the spine (spinal or epidural block) or around a major nerve, numbing a whole region of the body
    \item \textbf{Local anesthesia:} an anesthetic is applied or injected into the skin and tissues around a small site, such as for a dental filling or a skin stitch
    \item \textbf{Sedation:} lighter levels of medication are used for procedures where the person should remain relaxed but able to breathe on their own
    \item \textbf{Monitoring:} the anesthesiologist continually checks heart rate, blood pressure, breathing, and oxygen levels, and adjusts the depth of anesthesia accordingly
    \item \textbf{Recovery:} as the medicines wear off, consciousness, sensation, and muscle control gradually return
\end{itemize}

\section*{Clinical Features}
\begin{itemize}
    \item Complete pain relief during the procedure, and for a time afterwards
    \item With general anesthesia, loss of consciousness and of memory of the procedure
    \item With regional anesthesia, numbness and loss of movement in the numbed area, lasting for hours
    \item Common after-effects: drowsiness, a sore throat from the breathing tube, nausea, shivering, and dizziness
    \item Return of feeling, movement, and full awareness as the anesthetic wears off
    \item Signs of a problem that need attention: persistent nausea, breathing difficulty, or pain that is not controlled
\end{itemize}

\section*{Differential Diagnosis}
Because anesthesia is a planned treatment rather than a disease, the key distinctions involve choosing the right technique and interpreting symptoms afterwards:
\begin{itemize}
    \item General versus regional versus local anesthesia, chosen according to the procedure and the person
    \item Normal after-effects of anesthesia, such as nausea, drowsiness, or a sore throat, versus signs of a complication
    \item Surgical pain versus pain or numbness caused by the anesthesia itself, such as a nerve block that is slow to wear off
    \item Prolonged drowsiness versus a genuine medical event, such as low blood sugar or a neurological problem
\end{itemize}

\section*{When to Seek Medical Care}
Before any procedure, tell your anesthesiologist about your medical conditions, medicines, allergies, and any previous problems with anesthesia. After the procedure, seek medical advice if you are concerned about your recovery.

\textbf{Call 000 (or local emergency services) for:}
\begin{itemize}
    \item Difficulty breathing or chest pain after a procedure
    \item Fainting, collapse, or reduced consciousness
    \item A severe allergic-type reaction, such as sudden difficulty breathing or widespread hives
    \item Numbness, weakness, or pain that is severe, worsening, or not resolving as expected
\end{itemize}

\section*{Diagnosis and Investigations}
\begin{itemize}
    \item \textbf{Pre-operative assessment:} a review of health, medicines, allergies, and previous anesthesia experiences before the procedure
    \item \textbf{Physical examination:} including the heart and lungs, and assessment of the airway
    \item \textbf{Blood tests:} depending on the person and procedure, such as a full blood count or kidney function
    \item \textbf{Heart tests:} an ECG, and sometimes further cardiac assessment, if there is heart disease
    \item \textbf{Lung function:} checks such as spirometry for people with lung disease
    \item \textbf{Discussion of choices:} the anesthesiologist explains the options, risks, and what will happen, and obtains informed consent
\end{itemize}

\section*{Management}
\begin{itemize}
    \item \textbf{Choice of technique:} the safest and most suitable type of anesthesia is selected for each person and procedure
    \item \textbf{Preparation:} fasting before general or regional anesthesia to reduce the risk of inhaling stomach contents
    \item \textbf{Delivery and monitoring:} medicines are given and continuously monitored by the anesthesiologist throughout the procedure
    \item \textbf{Pain relief during and after:} combinations of medicines, nerve blocks, and local anesthetics control pain in the recovery period
    \item \textbf{Nausea management:} medicines to prevent and treat postoperative nausea and vomiting
    \item \textbf{Recovery care:} monitoring in the recovery room until awake and stable, then ongoing pain relief and early mobilization
\end{itemize}

\section*{Complications}
\begin{itemize}
    \item Nausea and vomiting after general anesthesia
    \item A sore throat, dry mouth, or hoarse voice from the breathing tube
    \item Drowsiness, confusion, or memory lapses, which may last longer in older people (postoperative delirium)
    \item Breathing problems, including a slow recovery of breathing or, rarely, inhaling stomach contents
    \item Allergic reactions to anesthetic medicines, which are uncommon
    \item Nerve damage or persistent pain or numbness at an injection site after regional anesthesia
    \item Heart and blood pressure changes, particularly in people with existing heart disease
\end{itemize}

\section*{Prognosis and Outcomes}
For most people, anesthesia is very safe and the after-effects settle within hours to a day. Nausea and drowsiness usually resolve quickly, and any sore throat lasts a day or two. Regional anesthesia typically wears off within several hours. Serious complications are rare in healthy people, and the risks are lower now than in the past because of better monitoring and medicines. Outcomes are best when people attend the pre-operative assessment honestly, follow fasting and medication instructions, and report any unusual symptoms afterwards.

\section*{Prevention and Self-Care}
\begin{itemize}
    \item Attend the pre-operative assessment and give a full, accurate list of medicines, allergies, and health conditions
    \item Follow fasting instructions exactly, usually with no food or drink for a set number of hours before anesthesia
    \item Ask questions until you understand the plan and what to expect during recovery
    \item Arrange transport home after procedures and avoid driving, operating machinery, or signing important documents until told it is safe
    \item Take prescribed pain relief as directed and mobilize gently when advised
    \item Report any persistent or unusual symptoms to your medical team after the procedure
\end{itemize}

\section*{Sources and Further Reading}
The following resources provide reliable, peer-reviewed and patient-orientated information on anesthesia.

\begin{itemize}
    \item NHS. \textit{Overview of anesthesia, including general and regional anesthesia.} https://www.nhs.uk/conditions/
    \item Mayo Clinic. \textit{Guide to the types of anesthesia and what to expect.} https://www.mayoclinic.org/
    \item MSD Manual. \textit{Professional and patient information on anesthesia and its risks.} https://www.msdmanuals.com/
    \item MedlinePlus. \textit{Health information on anesthesia and anesthesia medicines.} https://medlineplus.gov/
    \item NICE. \textit{Clinical guidance on perioperative care and patient safety.} https://www.nice.org.uk/
\end{itemize}
